% Corrected French for printed p.180; departures are append-only below.
\SourcePage{180}

\LayerTarget{tohoku:corollary:3.9.1:3}
\textsc{Corollaire 3.} \textit{Supposons la condition préliminaire du
corollaire précédent satisfaite, supposons de plus que les recouvrements
$R\in\mathfrak R(A)$ aient un nerf de dimension $\leqq n$. Alors les
conditions équivalentes du corollaire précédent impliquent même
$H^i(A,F)=0$ pour $A\in\mathfrak S$ et tout $i>0$.}

\LayerTarget{tohoku:proof:corollary-3.9.1-3}
On aura en effet automatiquement $H^i(R,F)=0$ pour
$\mathbf R\in\mathfrak R(A)$ et $i>n$, donc pour tout $i>0$, c'est-à-dire la
condition $A(m)$ sera vérifiée pour \emph{tout $m$}, donc il suffit d'appliquer le
corollaire avec $m$ grand.

\LayerTarget{tohoku:proof:theorem-3.9.1:induction-on-n}
\textsc{Démonstration du Théorème 3.9.1.} Nous procéderons par
récurrence sur l'entier $n$, la proposition étant évidente si $n=0$.
Supposons donc $n\geqq1$, et le théorème démontré pour les entiers $n'<n$. De
l'hypothèse de récurrence on tire d'abord aussitôt $H^i(A,F)=\{0\}$ pour
$1\leqq i\leqq n-1$, il reste donc à prouver la nullité de
$c^n\in H^n(A,F)$, sous la condition qu'il existe une partie finie $L$ de
$\mathfrak R(A)$ telle que la restriction de $c^n$ à tout
$B\in\mathbf R^L$ soit nulle. Soit $k$ le nombre d'éléments de $L$, nous
raisonnons par récurrence sur $k$. La conclusion est triviale si $k=0$,
prouvons-la pour $k=1$.
\LayerTarget{tohoku:proof:theorem-3.9.1:base-k-one}
Par hypothèse, il existe $\mathbf R\in\mathfrak R(A)$ tel que la restriction
de $c^n$ à tout $B\in\mathbf R$ soit nulle. Comme les $B\in\mathbf R$ sont
dans $\mathfrak S$, d'où $H^i(B,F)=\{0\}$ pour $1\leqq i\leqq n-1$, le terme
$E_2^{p,q}$ de la suite spectrale de Leray relative à $\mathbf R$ est nul pour
$1\leqq q\leqq n-1$. On en conclut classiquement une suite exacte
\LayerTarget{tohoku:eq:theorem-3.9.1-edge-sequence}
$H^n(\mathbf R,F)\to H^n(A,F)
\to H^0(\mathbf R,H^n(F))\to\ldots$.
Comme le premier terme est nul en vertu de $A(n)$, le deuxième homomorphisme
est injectif, or par hypothèse l'image de $c^n$ par ce dernier est nulle, donc
$c^n$ est nul.
\LayerTarget{tohoku:proof:theorem-3.9.1:induction-on-k}
Supposons maintenant $k\geqq2$, et la conclusion démontrée pour les valeurs
$k'<k$. Soit $L=(\mathbf R^1,\ldots\mathpunct{.},\mathbf R^k)$. Il suffit de prouver,
d'après ce qu'on a vu, que la restriction $c_B^n$ de $c^n$ à tout
$B\in R^1$ est nulle. Or pour $i=2,\ldots\mathpunct{.},k$, la restriction
$\mathbf R_B^i$ de $\mathbf R^i$ à $B$ appartient à $\mathfrak R(B)$,
d'autre part la restriction de $c_B^n$ à tout ensemble appartenant à
l'intersection des recouvrements $\mathbf R_B^i$ ($2\leqq i\leqq k$) est nulle
par hypothèse sur $c^n$. Appliquant notre hypothèse de récurrence, pour
$k'=k-1$, à $c_B^n$ et $B$, on trouve que $c_B^n=0$, ce qui achève la
démonstration du théorème.

\LayerTarget{tohoku:proposition:3.9.2}
\textsc{Proposition 3.9.2.} \textit{L'hypothèse préliminaire du
corollaire \textup{2} est satisfaite dans le cas suivant : $X$ est quasicompact, les
$A\in\mathfrak S$ sont fermés, $X\in\mathfrak S$, les
$\mathbf R\in\mathfrak R(X)$ sont finis et, pour tout $A\in\mathfrak S$ et
tout $\mathbf R\in\mathfrak R(X)$, la trace $\mathbf R_A$ de $\mathbf R$ sur
$A$ appartient à $\mathfrak R(A)$ ; de plus, pour deux points distincts $x,y$
de $X$, il existe un $\mathbf R\in\mathfrak R(X)$ dont aucun ensemble ne
contient à la fois $x$ et $y$.}

\LayerTarget{tohoku:definition:quasi-compact}
(Nous dirons qu'un espace est quasi-compact s'il satisfait à l'axiome des
recouvrements ouverts des espaces compacts, sans être néanmoins nécessairement
séparé).

\LayerTarget{tohoku:construction:cover-stars-and-neighborhoods}
\textsc{Démonstration.} Pour un recouvrement $\mathbf R$ de $X$ et un
$x\in X$, désignons par $E_x(\mathbf R)$ (\emph{étoile} de $\mathbf R$ en $x$) la
réunion des $A\in\mathbf R$ qui contiennent $x$, et soit $O_x(\mathbf R)$ le
complémentaire de la réunion des $A\in\mathbf R$ qui ne contiennent pas $x$.
\LayerTarget{tohoku:eq:cover-star-inclusions}
On a donc $x\in O_x(\mathbf R)\subset E_x(\mathbf R)$, de plus pour tout
$y\in O_x(\mathbf R)$, on a $E_y(\mathbf R)\subset E_x(\mathbf R)$.
\LayerTarget{tohoku:claim:closed-cover-star-neighborhood}
Si $\mathbf R$ est un recouvrement fini de $X$ par des ensembles fermés,
$O_x(\mathbf R)$ est un \textit{voisinage ouvert} de $x$.
\LayerTarget{tohoku:proof:proposition-3.9.2:finite-intersection:p1}
Sous les conditions de la proposition, soit $\mathbf U$ un recouvrement
ouvert de $X$, soit $x\in X$, soit $U_x\in\mathbf U$ avec $x\in U_x$ ;
l'intersection

\TohokuContinuation{181}
